\clearpage
\markboth{作者简介}{作者简介}
\phantomsection
\addcontentsline{toc}{section}{作者简介}
\begingroup
    \centering
    \zihao{-3}\heiti\bfseries 作者简介
    \par
    \vspace{\baselineskip}
\endgroup

\begingroup
    \zihao{-4}\songti
    \setlength{\baselineskip}{\bodylineskip}
    \parindent 2em
    \textbf{XXX}，男（XXXX--），1994 年毕业于 XXXX，获学士学位；XXXX 年毕业于 XXXX，获硕士学位；XXXX 年 X 月--XXXX 年 X 月在中国矿业大学（北京）攻读博士学位，专业为 XXXX，研究方向为 XXXX。
\endgroup

\vspace{2\baselineskip}
\begingroup
    \centering
    \zihao{4}\heiti\bfseries 在学期间发表学术论文
    \par
    \vspace{\baselineskip}
\endgroup

\begin{enumerate}[label=\arabic*., leftmargin=2em, itemsep=0pt, parsep=0pt]
    \zihao{-4}\songti
    \item \textbf{XXX}, Ren Deyi, Yang Jianye et al. TOF-SIMS study of the hydrocarbon-generating potential of mineral-bituminous groundmass[J]. Acta Geological Sinica, 2000, 74(1): 203-208. (SCI 收录)(学位论文第*章)
    \item \textbf{XXX}. The action and significance of low organism in the formation of high-sulfur coal[J]. Scientia Geologica Sinica, 2000, 9(3): 222-229. (学位论文第*章)
    \item \textbf{XXX}，任德贻，唐跃刚等．高硫煤中硫的地质演化模式——以内蒙古乌达矿区为例[J]．地质论评，2001，47(4)：16-25.（学位论文第*章）
    \item \textbf{XXX}，任德贻，艾天杰等．乌达矿区主采煤层可选性评价的地质因素研究[J]．中国矿业大学学报，2000，29(3)：305-312. (EI 收录)(学位论文第*章)
    \item \textbf{XXX}，艾天杰，焦方立等．内蒙古乌达矿区煤中硫的同位素组成及演化特征[J]．岩石学报，2000，16(2)：25-35. (SCI 收录)(学位论文第*章)
    \item \textbf{XXX}，唐跃刚，杨建业等．烃源岩生烃性的飞行时间二次离子质谱研究[J]．中国矿业大学学报，2000，29(6)：405-410. (EI 收录)(学位论文第*章)
\end{enumerate}

\vspace{1\baselineskip}
\begingroup
    \centering
    \zihao{4}\heiti\bfseries 在学期间出版学术专著
    \par
    \vspace{\baselineskip}
\endgroup
\begin{enumerate}[label=\arabic*., leftmargin=2em, itemsep=0pt, parsep=0pt]
    \zihao{-4}\songti
    \item \textbf{XXX}．生物有机质成矿作用和成矿背景[M]．北京：海洋出版社，1998.
\end{enumerate}

\vspace{1\baselineskip}
\begingroup
    \centering
    \zihao{4}\heiti\bfseries 在学期间授权专利
    \par
    \vspace{\baselineskip}
\endgroup
\begin{enumerate}[label=\arabic*., leftmargin=2em, itemsep=0pt, parsep=0pt]
    \zihao{-4}\songti
    \item \textbf{XXX}．全智能节电器：200610171314.3[P]．2006-12-13.
\end{enumerate}

\clearpage

\begingroup
    \centering
    \zihao{4}\heiti\bfseries 在学期间主要获奖
    \par
    \vspace{\baselineskip}
\endgroup
\begin{enumerate}[label=\arabic*., leftmargin=2em, itemsep=0pt, parsep=0pt]
    \zihao{-4}\songti
    \item “煤的洁净利用地质技术”于 2001 年获得教育部中国高校科技进步一等奖（排名第 2）。
    \item 2001 年获第十三届邝寿堃奖学金。
\end{enumerate}

\vspace{1\baselineskip}
\begingroup
    \centering
    \zihao{4}\heiti\bfseries 在学期间参与科研项目
    \par
    \vspace{\baselineskip}
\endgroup
\begin{enumerate}[label=\arabic*., leftmargin=2em, itemsep=0pt, parsep=0pt]
    \zihao{-4}\songti
    \item 国家自然科学基金项目“煤中有害元素富集成因类型-以鄂尔多斯晚古生代煤为例”主要研究人员。项目编号：40072054，2001 年 1 月-2003 年 12 月。
    \item 煤炭科学基金项目“煤中硫的污染抑制性与可选性评价的地质成因”课题负责人。项目编号：97 地 10205，1998 年 1 月-2000 年 12 月。
    \item 国家自然科学基金项目“沥青质体和矿物沥青基质的来源，形成期次及生烃性研究”主要研究人员。编号：49772132，1998 年 1 月-2000 年 12 月。
\end{enumerate}
